\section{A Liouville--cyclotomic product identity}\label{sec:liouville}

Let $\lambda(n)=(-1)^{\Omega(n)}$ be the Liouville function, where
$\Omega(n)$ counts prime factors with multiplicity, and let $\Phi_n$ be
the $n$th cyclotomic polynomial. Stuart Clary's sequence A118207 is defined by
the formal generating function
\begin{equation}\label{eq:liouville-A}
 A(z)=\prod_{n\ge1}(1+z^n)^{\lambda(n)}.
\end{equation}
Peter Bala's formula entry of 5 April 2023 conjectures the following
identity \cite{oeis118207}.

\begin{theorem}\label{thm:liouville-main}
In the ring of rational formal power series,
\begin{equation}\label{eq:liouville-target}
 A(z^2)=\prod_{n\ge1}\Phi_n(z^n)\Phi_n((-z)^n).
\end{equation}
\end{theorem}

Every product below is interpreted coefficientwise. In particular, its
coefficient in degree $N$ is determined by the factors with index at most
$\max\{1,N\}$. Integer exponents are powers in the group of series with constant
coefficient $1$. At index $1$ the individual cyclotomic factors have
constant coefficient $-1$, but their paired product has constant
coefficient $1$.

The proof uses a product identity already established in Proposition 1 of
Bala's 9 April 2023 note \cite{bala2023}: if
\[
 F(z)=\prod_{n\ge1}(1-z^n)^{\lambda(n)},
\]
then
\begin{equation}\label{eq:bala-known}
 F(z)=-\prod_{n\ge1}\Phi_n(z^n).
\end{equation}
The inspected four-page note proves this identity and a companion
formula for its reciprocal. It does not state
\eqref{eq:liouville-target}. We include a proof of the helper identity to
make the argument and its formalization self-contained, while retaining
Bala's attribution.

\begin{lemma}[Square-divisor identity]\label{lem:liouville-square}
For every positive integer $m$,
\begin{equation}\label{eq:liouville-square}
 \lambda(m)=\sum_{d^2\mid m}\mu(m/d^2),
\end{equation}
where $\mu$ denotes the M\"obius function. Also,
$\lambda(2m)=-\lambda(m)$.
\end{lemma}

\begin{proof}
Write $m=\prod_p p^{e_p}$. There is a unique factorization
$m=ab^2$ with $a$ squarefree: the exponent of $p$ in $a$ is
$e_p\bmod2$, and its exponent in $b$ is $\lfloor e_p/2\rfloor$.
Every term in the sum is zero unless $m/d^2$ is squarefree, so the unique
nonzero term is
\[
 \mu(a)=(-1)^{\sum_p(e_p\bmod2)}
        =(-1)^{\sum_p e_p}=\lambda(m).
\]
Complete multiplicativity follows by adding prime exponents, and
$\lambda(2)=-1$ gives the final assertion.
\end{proof}

\begin{lemma}[Bala's normalized Euler product]\label{lem:liouville-bala}
Define $\Psi_1(z)=1-z$ and $\Psi_n(z)=\Phi_n(z)$ for $n>1$.
Then
\begin{equation}\label{eq:liouville-normalized}
 \prod_{n\ge1}\Psi_n(z^n)=F(z).
\end{equation}
\end{lemma}

\begin{proof}
The usual cyclotomic factorization, with the sign at $n=1$ normalized,
gives
\[
 \prod_{d\mid n}\Psi_d(z)=1-z^n.
\]
M\"obius inversion in the multiplicative group of formal units gives
\begin{equation}\label{eq:liouville-mobius}
 \Psi_n(z)=\prod_{d\mid n}(1-z^d)^{\mu(n/d)}.
\end{equation}
After substituting $z^n$ and writing $n=ad$, the product becomes
\[
 \prod_{a,d\ge1}(1-z^{ad^2})^{\mu(a)}.
\]
Regroup by $m=ad^2$. The exponent of $1-z^m$ is
$\sum_{d^2\mid m}\mu(m/d^2)=\lambda(m)$ by
Lemma~\ref{lem:liouville-square}. This proves
\eqref{eq:liouville-normalized}. For a fixed degree $N$, only pairs with
$ad^2\le N$ can contribute, so the regrouping is finite modulo
$z^{N+1}$. The single sign change at index $1$ recovers
\eqref{eq:bala-known}.
\end{proof}

\begin{proof}[Proof of Theorem~\ref{thm:liouville-main}]
First compute $F(z)F(-z)$. For odd $n$,
$(1-z^n)(1-(-z)^n)=1-z^{2n}$; for even $n$ the product is
$(1-z^n)^2$. Splitting odd and even indices, then writing $n=2m$ in
the even part, gives
\begin{align}
 F(z)F(-z)
 &=\prod_{\substack{m\ge1\\m\text{ odd}}}
       (1-z^{2m})^{\lambda(m)}
   \prod_{m\ge1}(1-z^{2m})^{2\lambda(2m)}\notag\\
 &=\prod_{\substack{m\ge1\\m\text{ odd}}}
       (1-z^{2m})^{\lambda(m)}
   \prod_{m\ge1}(1-z^{2m})^{-2\lambda(m)}.
   \label{eq:liouville-evenodd}
\end{align}
Thus the exponent attached to $1-z^{2m}$ is $-\lambda(m)$ when
$m$ is odd and $-2\lambda(m)$ when $m$ is even.

On the other hand, factoring $1-z^{4n}$ gives
\begin{align*}
 A(z^2)
 &=\prod_{n\ge1}
       \left(\frac{1-z^{4n}}{1-z^{2n}}\right)^{\lambda(n)}\\
 &=\prod_{\substack{m\ge1\\m\text{ even}}}
       (1-z^{2m})^{-\lambda(m)}
   \prod_{m\ge1}(1-z^{2m})^{-\lambda(m)}.
\end{align*}
For the numerator reindexing we used $m=2n$ and
$\lambda(n)=-\lambda(m)$. The exponent at every $m$ therefore
agrees with \eqref{eq:liouville-evenodd}, proving
$A(z^2)=F(z)F(-z)$. Apply \eqref{eq:bala-known} at $z$ and $-z$.
The two overall minus signs cancel, yielding
\eqref{eq:liouville-target}.
\end{proof}

\paragraph{Formal products and the exact source bridge.}
The Lean proof first establishes these identities for an arbitrary
nilpotent element $x$ of any commutative ring. Each $1-x^n$ with $n>0$
is a unit, and all sufficiently high powers of $x$ are zero. Consequently
the products have finite multiplicative support. The proof of
Lemma~\ref{lem:liouville-bala} uses two explicit finite fiber regroupings
of the factors $(1-x^{ad^2})^{\mu(a)}$: first by $ad$, then by $ad^2$.
The square-divisor identity is proved arithmetically from the unique
squarefree-times-square decomposition, including uniqueness and the
vanishing of every other M\"obius term.

For each $N\ge0$ this nilpotent theorem is applied to the class of $X$ in
\[
 \mathbb{Q}[[X]]/(X^{N+1}).
\]
The literal factor families in \path{SourceSeries.lean} are
$(1+X^{2n})^{\lambda(n)}$ and
$\Phi_n(X^n)\Phi_n((-X)^n)$, with index zero set to the neutral factor.
The file proves that their quotient images equal $1$ for $n>N$ and
identifies the remaining images with the finite-support unit products.
The generic construction \path{formalProduct} takes the degree-$N$
coefficient from the prefix through index $N$. Its stabilization theorem
proves that this definition agrees with every sufficiently long prefix
in the quotient, rather than assuming a convergence property.

The final theorem \path{A118207.source_conjecture} proves equality of
the two resulting formal power series.
\path{A118207.source_coefficient_identity} states their equality in
every degree. Thus the formal conclusion is the complete infinite
product identity, with no fixed truncation bound. The methods and
the helper identity are classical; this package resolves the specific
conjectural formula recorded in A118207 and supplies its complete Lean
proof. The bounded prior-work search is documented separately and does
not establish worldwide priority.
